\section{Least discrepancy: from mismatch accumulation to phase potential and the Newtonian limit}
\label{sec:least_discrepancy}

\subsection{Discrepancy as a readout mismatch measure}
\label{subsec:discrepancy_definition}

For a point set $\{x_n\}_{n=1}^N\subset[0,1)$, define the one-dimensional \emph{star discrepancy} \cite{KuipersNiederreiter1974}
\begin{equation}
  D_N^{\ast} := \sup_{a\in[0,1]}\left|\frac{1}{N}\sum_{n=1}^N \ind_{[0,a)}(x_n) - a\right|,
  \qquad
  E_N := N D_N^{\ast}.
\end{equation}
In the scan model $x_n=x_0+n\alpha\pmod{1}$, $E_N$ is interpreted as the accumulated mismatch between finite-prefix readout statistics and the uniform ontic measure. CAP treats mismatch as \emph{structural}: it cannot be eliminated by longer observation but can be controlled.

To connect $D_N^\ast$ to scan readout stability, write the orbit as $x_n=x_0+n\alpha\pmod 1$ and consider interval-window counts
\begin{equation}
  s_n = \ind_W(x_0+n\alpha),
  \qquad
  S_N(W)=\sum_{n=0}^{N-1}s_n,
  \qquad
  W\subset[0,1)\ \text{an interval}.
\end{equation}
For $W=[0,a)$ one has $S_N(W)=\sum_{n=1}^N \ind_{[0,a)}(x_n)$ and therefore
\begin{equation}
  E_N = \sup_{a\in[0,1]}\left|S_N([0,a)) - Na\right|.
\end{equation}
Thus, a uniform interval-count bound immediately yields a star-discrepancy bound.

\paragraph{Continued fractions and Ostrowski expansion.}
Let $\alpha=[0;a_1,a_2,\dots]\in(0,1)\setminus\QQ$ and let $(q_k)_{k\ge 0}$ be the convergent denominators defined by
\begin{equation}
  q_{-1}=0,\quad q_0=1,\quad q_{k+1}=a_{k+1}q_k+q_{k-1}.
\end{equation}
Every $N\in\NN_{>0}$ admits a unique Ostrowski expansion relative to $\alpha$,
\begin{equation}
  N=\sum_{k=0}^m b_k q_k,
\end{equation}
with digits satisfying $0\le b_0<a_1$, $0\le b_k\le a_{k+1}$ for $k\ge 1$ and the standard admissibility constraint (if $b_k=a_{k+1}$ then $b_{k-1}=0$); see e.g.\ \cite{AlloucheShallit2003,Lothaire2002}.

\begin{theorem}[Ostrowski--Denjoy--Koksma bound for interval counts]
\label{thm:ostrowski_dk_cap}
Let $\alpha=[0;a_1,a_2,\dots]\in(0,1)\setminus\QQ$ and let $q_k$ be the convergent denominators.
Let $W\subset[0,1)$ be any interval and let $S_N(W)=\sum_{n=0}^{N-1}\ind_W(x_0+n\alpha)$.
For every $N\in\NN_{>0}$ write the Ostrowski expansion $N=\sum_{k=0}^m b_k q_k$.
Then for every $x_0\in[0,1)$,
\begin{equation}
  \left|S_N(W) - N\,\mu(W)\right|
  \le 2\sum_{k=0}^m b_k
  \le 2\sum_{k=0}^m a_{k+1},
\end{equation}
where $\mu$ is Lebesgue measure on $[0,1)$.
\end{theorem}

\begin{proof}
Let $f:=\ind_W-\mu(W)$. Then $f$ has bounded variation $\mathrm{Var}(f)=2$ and $\int_0^1 f\,\dd\mu=0$.
For each $k\ge 0$, the Denjoy--Koksma inequality for rotations gives a uniform bound at convergent times:
\begin{equation}
  \left|\sum_{j=0}^{q_k-1} f(x_0+j\alpha)\right|\le \mathrm{Var}(f)=2,
\end{equation}
see e.g.\ \cite{Denjoy1932,Koksma1935,KatokHasselblatt1995}.
Now expand $N=\sum_{k=0}^m b_k q_k$ and decompose the length-$N$ sum into $b_m$ blocks of length $q_m$, then $b_{m-1}$ blocks of length $q_{m-1}$, and so on (the standard Ostrowski block decomposition).
Applying the $q_k$-time bound to each block yields
\begin{equation}
  \left|\sum_{n=0}^{N-1} f(x_0+n\alpha)\right|
  \le 2\sum_{k=0}^m b_k.
\end{equation}
Since $\sum_{n=0}^{N-1} f(x_0+n\alpha)=S_N(W)-N\mu(W)$, the first inequality follows.
The second inequality follows from the digit bounds $b_0<a_1$ and $b_k\le a_{k+1}$ for $k\ge 1$.
\end{proof}

\begin{corollary}[Golden optimality for a discrepancy proxy]
\label{cor:golden_proxy_cap}
Fix a depth $m\ge 0$ and define the continued-fraction proxy
\begin{equation}
  C_m(\alpha):=\sum_{k=0}^m a_{k+1},
  \qquad \alpha=[0;a_1,a_2,\dots].
\end{equation}
Then $C_m(\alpha)\ge m+1$ for every irrational $\alpha$, with equality if and only if $a_i=1$ for all $1\le i\le m+1$ (the golden branch prefix).
In particular, the golden slope $\alpha=\varphi^{-1}=[0;1,1,1,\dots]$ uniquely minimizes the upper bound in Theorem~\ref{thm:ostrowski_dk_cap} at every finite depth.
\end{corollary}

\begin{proof}
Since each $a_i\in\NN$, one has $a_i\ge 1$ and therefore $C_m(\alpha)\ge m+1$, with equality if and only if $a_i=1$ for $1\le i\le m+1$.
The uniqueness of continued fractions yields the last statement.
\end{proof}

\begin{remark}[Star discrepancy bound and logarithmic growth for bounded type]
Applying Theorem~\ref{thm:ostrowski_dk_cap} to $W=[0,a)$ and taking the supremum over $a$ yields
\begin{equation}
  E_N = N D_N^\ast \le 2\sum_{k=0}^m a_{k+1}.
\end{equation}
If $\alpha$ is of bounded type, i.e.\ $\sup_k a_k\le A<\infty$, then $E_N\le 2A(m+1)$.
Moreover, since $a_{k+1}\ge 1$ implies $q_{k+1}\ge q_k+q_{k-1}$, one has $q_k\ge F_k$ (Fibonacci), hence $m=O(\log N)$ whenever $q_m\le N<q_{m+1}$. Therefore $E_N=O(\log N)$ for bounded type slopes.
\end{remark}

This quantitative control should be contrasted with resonance-prone slopes that admit exceptionally good rational approximations on intermediate scales. In CAP semantics, good rational approximation corresponds to \emph{phase locking} over finite depth: the scan visits window boundaries in structured patterns that amplify readout mismatch before equidistribution asserts itself. Least discrepancy is therefore not only a statement about asymptotic equidistribution but a statement about finite-resolution stability under repeated projection.

\subsection{Finite-depth optimality of the golden branch}
\label{subsec:golden_optimality}

The previous subsection provides a concrete finite-depth proxy $C_m(\alpha)=\sum_{k=0}^m a_{k+1}$ controlling mismatch via Denjoy--Koksma and Ostrowski decomposition. The golden slope
\begin{equation}
  \alpha_{\varphi} := \varphi^{-1} = \frac{\sqrt{5}-1}{2},
  \qquad
  [0;1,1,1,\dots],
\end{equation}
minimizes this proxy uniformly at every depth (Corollary~\ref{cor:golden_proxy_cap}) and is the archetype of ``most badly approximable'' numbers \cite{Cassels1957}.

\paragraph{Canonical ticks via Ostrowski truncation.}
At finite depth $m$, an observer effectively accesses only an $m$-truncated Ostrowski expansion of scan indices; this induces a discrete tick structure whose stability depends on the continued-fraction prefix \cite{AlloucheShallit2003,Lothaire2002}. In the golden branch, the truncation specializes to Zeckendorf/Fibonacci addressing, supplying a canonical ``clock'' without external tuning \cite{Zeckendorf1972}.

\noindent In CAP terms, the golden branch is not aesthetic but variational: for a fixed readout resolution, it minimizes the \emph{worst-case} upper bound controlling mismatch accumulation, thereby lowering the cost required for long-term self-consistent readout. Pointwise $E_N$ at a fixed $N$ and $x_0$ can still fluctuate among bounded-type slopes (Section~\ref{sec:numerics}); CAP's claim is the uniform stability of the bound that governs sustainable readout.

\subsection{Phase potential and phase pressure: mismatch sources generate $1/r$ fields}
\label{subsec:phase_potential}

Omega Dynamics lifts coarse-grained mismatch into a continuum source \cite{Ma2025HPD}. Let $\sigma(x)$ be a mismatch density obtained from a regulated continuum limit of $E_N$ (Convention R1). We define the \emph{phase potential} $\Phi$ as the stationary point of the quadratic functional $S_\Phi[\Phi;\sigma]$ in \eqref{eq:phase_energy_functional}; equivalently, $\Phi$ solves the Poisson equation
\begin{equation}
  \Delta \Phi = 4\pi \rho_{\Phi},
  \qquad
  \rho_{\Phi} = \kappa_{\Phi}\, \sigma .
\end{equation}
We then define the \emph{phase pressure} vector field
\begin{equation}
  \mathbf{P}_{\Phi} := -\nabla \Phi .
\end{equation}
The Poisson form is the natural weak-field, slow-variation limit: in the Newtonian limit of GR one obtains $\Delta\phi=4\pi G\rho$ for the gravitational potential $\phi$ \cite{Wald1984,Carroll2004}. CAP uses the same operator as the minimal continuum lift of mismatch accumulation when promoted to an effective potential.

\paragraph{Variational closure for the phase potential.}
Independently of GR, the Poisson equation is the Euler--Lagrange equation of the quadratic field functional
\begin{equation}
\label{eq:phase_energy_functional}
  S_{\Phi}[\Phi;\sigma]
  := \int_{\RR^{3}}\dd^{3}x\,
  \left(
    \frac{1}{8\pi}\,|\nabla\Phi|^{2}
    + \kappa_{\Phi}\,\sigma\,\Phi
  \right),
\end{equation}
with fixed source $\sigma$ and appropriate decay/boundary conditions. Varying $\Phi\mapsto\Phi+\epsilon \delta\Phi$ and integrating by parts yields
\begin{equation}
  \delta S_{\Phi}
  = \int \dd^{3}x\,\delta\Phi\left(-\frac{1}{4\pi}\Delta\Phi+\kappa_{\Phi}\sigma\right),
  \qquad\Rightarrow\qquad
  \Delta\Phi = 4\pi\kappa_{\Phi}\sigma.
\end{equation}
This recovers the Poisson-type sourcing used above. CAP therefore treats the Poisson lift not as an extra postulate but as the \emph{least-action} closure for a scalar potential field that encodes mismatch cost at the macroscopic level.

\paragraph{Exterior solution and mismatch charge.}
Let $\sigma$ be compactly supported in $\RR^3$ and assume $\Phi\to 0$ at spatial infinity. Using the Green function identity $\Delta(1/|\mathbf{x}|)=-4\pi\delta(\mathbf{x})$ \cite{ArfkenWeberHarris2013}, the unique decaying solution of
\begin{equation}
  \Delta\Phi = 4\pi \kappa_\Phi \sigma
\end{equation}
is
\begin{equation}
  \Phi(\mathbf{x})
  = -\kappa_\Phi\int_{\RR^3}\frac{\sigma(\mathbf{y})}{|\mathbf{x}-\mathbf{y}|}\,\dd^3\mathbf{y}.
\end{equation}
In particular, for $r=|\mathbf{x}|\to\infty$ one has the multipole expansion
\begin{equation}
  \Phi(r)
  = -\frac{M}{r} + O(r^{-2}),
  \qquad
  M := \kappa_\Phi \int_{\RR^3}\sigma(\mathbf{y})\,\dd^3\mathbf{y}.
\end{equation}
Thus the monopole coefficient $M$ is fixed by the total mismatch charge $\int\sigma$ and the calibration constant $\kappa_\Phi$.

\begin{remark}[Signed discrepancy vs.\ mismatch charge]
The empirical measure $\mu_N=\frac1N\sum_{n=1}^N \delta_{x_n}$ and the uniform measure $\mu$ satisfy $\int(\mu_N-\mu)=0$, so a \emph{signed} equidistribution defect would have zero total charge and would not generate a monopole $1/r$ term. CAP's $\sigma$ is therefore not the signed measure $\mu_N-\mu$; it is a \emph{nonnegative coarse-grained cost density} associated with sustaining readout consistency (and, in the Newtonian closure of Appendix~\ref{app:newtonian_limit}, it can be tied to an effective energy density). This is precisely why an isolated defect sector can carry nonzero total mismatch charge and produce a monopole potential.
\end{remark}

\begin{remark}[Flux characterization]
By the divergence theorem and $\Delta\Phi=4\pi\kappa_\Phi\sigma$, one has
\begin{equation}
  \int_{S_R}\mathbf{P}_\Phi\cdot \dd\mathbf{S}
  = \int_{B_R}\nabla\cdot \mathbf{P}_\Phi\,\dd^3\mathbf{x}
  = -\int_{B_R}\Delta\Phi\,\dd^3\mathbf{x}
  = -4\pi \kappa_\Phi \int_{B_R}\sigma\,\dd^3\mathbf{x},
\end{equation}
so in the limit $R\to\infty$ the monopole strength satisfies
\begin{equation}
  M = -\frac{1}{4\pi}\lim_{R\to\infty}\int_{S_R}\mathbf{P}_\Phi\cdot \dd\mathbf{S}.
\end{equation}
This matches the topological-flux perspective in Section~\ref{sec:matter_defects}.
\end{remark}

For an isolated defect source with quantized flux (a topological input analogous to Dirac quantization; Section~\ref{sec:matter_defects}) \cite{Dirac1931,WuYang1975}, the exterior solution is necessarily harmonic and (under standard isolation assumptions) spherically symmetric, hence
\begin{equation}
  \Phi(r) = -\frac{M}{r},
  \qquad
  \mathbf{a}(r) \equiv \mathbf{P}_{\Phi}(r) = -\frac{M}{r^{2}}\, \hat{\mathbf{r}} ,
\end{equation}
recovering the Newtonian inverse-square form as a mismatch-induced phase-pressure limit.

\noindent This establishes CAP's first closed chain:
\begin{equation}
  \text{discrete readout mismatch (discrepancy)}
  \;\Rightarrow\;
  \text{phase potential } \Phi
  \;\Rightarrow\;
  \text{phase pressure}
  \;\Rightarrow\;
  \text{Newtonian limit}.
\end{equation}


